On considère les couples acido-basiques~:

\begin{itemize}
    \item acide sulfurique/ion hydrogénosulfate~:
          $\ce{H2SO4}\ttslash{}\ce{HSO^-_4_{(aq)}}$

    \item ion hydrogénosulfate/ion sulfate~:
          $\ce{HSO^-_4_{(aq)}}\ttslash{}\ce{SO^2-_4_{(aq)}}$
\end{itemize}


\begin{enumerate}
    \item Quels rôles l'ion hydrogénosulfate joue-t-il dans ces deux couples~?
          Comment le qualifie-t-on~?
    \item Quel est le caractère acido-basique de l'ion $\ce{HO^-_{(aq)}}$~? Dans
          quelle solution le trouve-t-on~? Quelle est son espèce conjuguée~?
    \item L'ion $\ce{HO^-_{(aq)}}$ peut-il réagir avec l'acide sulfurique~?
          Si oui, écrire l'équation de la réaction acido-basique (étape~1).
    \item L'ion $\ce{HO^-_{(aq)}}$ peut-il réagir avec l'ion hydrogénosulfate~?
          Si oui, écrire l'équation de la réaction acidobasique (étape~2).
    \item En combinant les deux étapes, écrire l'équation de la réaction entre
          l'acide sulfurique et les ions $\ce{HO^-_{(aq)}}$.
          On dit que l'acide sulfurique est un diacide. Pourquoi~?
\end{enumerate}
